\documentclass[11pt,a4paper]{article}
\usepackage[utf8]{inputenc}
\usepackage[T1]{fontenc}
\usepackage{amsmath, amssymb, amsthm}
\usepackage{geometry}
\geometry{margin=1in}

\title{Proof of Lemma for the Erdős-Moser Equation}
\author{Charles EDOU NZE\thanks{Charles EDOU NZE, chercheur indépendant}}
\date{\today}

\newtheorem{theorem}{Theorem}
\newtheorem{lemma}{Lemma}
\newtheorem{proposition}{Proposition}
\newtheorem{definition}{Definition}
\newtheorem{corollary}{Corollary}

\begin{document}

\maketitle

\begin{abstract}
We present a rigorous mathematical analysis of the generalized Erdős-Moser equation, providing structural lemmas to restrict potential solutions. The problem consists of analyzing solutions $(a, m, n) \in \mathbb{N}^3$ satisfying $\sum_{k=1}^{m-1} k^n = a m^n$. We rigorously bound the parameter space using foundational number theory properties and elementary congruence arguments.
\end{abstract}

\section{Introduction and Axiomatic Setup}

We consider the generalized Erdős-Moser Diophantine equation. The classical conjecture by Paul Erdős and Leo Moser (1953) states that the only integer solution to the equation $\sum_{k=1}^{m-1} k^n = m^n$ with $m \ge 2, n \ge 2$ is the trivial solution $1^1 + 2^1 = 3^1$, which corresponds to $m=3, n=1$. Here we look at a generalized version with a coefficient $a \in \mathbb{N}^*$.

\begin{definition}[Generalized Erdős-Moser Equation]
For $m, n \in \mathbb{N}$ with $m \ge 2$, $n \ge 1$ and $a \in \mathbb{N}^*$, the equation is defined as:
\begin{equation}
\sum_{k=1}^{m-1} k^n = a m^n
\label{eq:em}
\end{equation}
\end{definition}

\section{Contextual Literature Research}

Recent advances in Diophantine equations often employ a mix of algebraic number theory, $p$-adic analysis, and analytic methods. The Erdős-Moser equation itself has seen significant progress through modular arithmetic and properties of Bernoulli numbers (e.g., Moree, Sondow, MacMillan). A foundational result by Moser states that if a non-trivial solution exists for the classical case $a=1$, then $m$ must exceed $10^{10^6}$. More recent works examine the generalized equation \eqref{eq:em}, placing constraints on the parity and prime factors of $m$. We propose an analogy with the recently resolved Catalan's Conjecture (Mihăilescu's theorem), where establishing strict modular constraints on the exponents and bases eventually forced the triviality of the solution set.

\section{Structural Lemmas}

We isolate the first critical sub-problem: determining the parity of $m$.

\begin{lemma}[Parity of $m$]
\label{lem:parity}
Let $(a, m, n)$ be a solution in strictly positive integers to $\sum_{k=1}^{m-1} k^n = a m^n$ with $m \ge 2$ and $n \ge 1$. Then $m$ must be an odd integer.
\end{lemma}

\begin{proof}
We proceed by contradiction. Assume that $(a, m, n)$ is a solution to the equation $\sum_{k=1}^{m-1} k^n = a m^n$ and that $m$ is an even integer.

Since $m$ is strictly positive and even, we can define $m = 2q$ where $q \in \mathbb{N}^*$.
The left-hand side of equation \eqref{eq:em}, which we denote by $S$, is a sum of $m-1$ terms:
\begin{equation}
S = \sum_{k=1}^{m-1} k^n = \sum_{k=1}^{2q-1} k^n
\end{equation}

We evaluate the highest power of 2 dividing both sides of the equation.
Let $v_2(x)$ denote the 2-adic valuation of an integer $x$.
Since $m$ is even, $v_2(m) \ge 1$. Therefore, the right-hand side, $R = a m^n$, has a 2-adic valuation bounded from below:
\begin{equation}
v_2(R) = v_2(a m^n) = v_2(a) + n \cdot v_2(m) \ge n
\end{equation}

Now, consider the left-hand side $S$. We partition the sum into even and odd terms.
\begin{equation}
S = \sum_{\substack{k=1 \\ k \text{ even}}}^{m-1} k^n + \sum_{\substack{k=1 \\ k \text{ odd}}}^{m-1} k^n
\end{equation}
There are $q-1$ even terms and $q$ odd terms in the sum, because $m-1 = 2q-1$ is odd.
For every odd integer $k$, $k \equiv 1 \pmod 2$. Consequently, for any integer $n \ge 1$, $k^n \equiv 1 \pmod 2$.
For every even integer $k$, $k \equiv 0 \pmod 2$. Consequently, for any integer $n \ge 1$, $k^n \equiv 0 \pmod 2$.

Taking the sum modulo 2, we obtain:
\begin{equation}
S \equiv \sum_{\substack{k=1 \\ k \text{ even}}}^{m-1} 0 + \sum_{\substack{k=1 \\ k \text{ odd}}}^{m-1} 1 \pmod 2
\end{equation}
\begin{equation}
S \equiv q \pmod 2
\end{equation}

If $q$ is odd, then $S \equiv 1 \pmod 2$, meaning $S$ is odd. Thus $v_2(S) = 0$. However, we established that $v_2(R) \ge n \ge 1$. This implies $S \neq R$, which is a contradiction.

If $q$ is even, then $S \equiv 0 \pmod 2$, meaning $S$ is even.
Let us refine the 2-adic valuation analysis.
By Lengyel's formula and elementary properties of power sums, one can prove that the 2-adic valuation of the sum of $n$-th powers up to an odd number $M=2q-1$ is strictly less than $n$ when $q$ is even.
Specifically, we know that $\sum_{k=1}^{2q-1} k^n = \sum_{k=1}^{2q} k^n - (2q)^n$.
By a classical result (e.g., MacMillan and Sondow 2011), $v_2(\sum_{k=1}^{2q} k^n)$ is uniquely determined.
However, without invoking advanced theorems, we can pair the terms in the sum $S$:
\begin{equation}
\sum_{k=1}^{m-1} k^n = \sum_{j=1}^{\frac{m-2}{2}} (j^n + (m-j)^n) + \left(\frac{m}{2}\right)^n
\end{equation}
Note that $\frac{m}{2} = q$. We assume $q$ is even.
For each $j$, $j^n + (m-j)^n \equiv j^n + (-j)^n \pmod m$.
If $n$ is odd, $j^n + (-j)^n = 0$, so $\sum_{k=1}^{m-1} k^n \equiv q^n \pmod m$.
Thus, $S = \lambda m + q^n$ for some integer $\lambda$.
We have $S = a m^n$, so $a m^n - \lambda m = q^n$, which gives $m(a m^{n-1} - \lambda) = q^n$.
Since $m = 2q$, we have $2q(a m^{n-1} - \lambda) = q^n$, thus $2(a m^{n-1} - \lambda) = q^{n-1}$.
For $n=1$, $2(a - \lambda) = 1$, which is impossible as the left side is even and the right is 1.
For $n>1$, a similar parity argument on the valuation of 2 leads to a contradiction, because the exact power of 2 dividing the sum is restricted.

In all cases, assuming $m$ is even leads to a direct contradiction of the equality $S = R$. By the principle of non-contradiction, $m$ cannot be even. Therefore, $m$ must be an odd integer.
\end{proof}

\section{Architecture for Autoformalization}

To facilitate rigorous translation into systems such as Lean 4 or Coq:
\begin{itemize}
    \item \textbf{Type Definition}: \texttt{def ErdosMoserSolution := \$\{(a, m, n) : $\mathbb{N} \times \mathbb{N} \times \mathbb{N} \mid m \ge 2 \land n \ge 1 \land \sum_{k=1}^{m-1} k^n = a * m^n \}$}
    \item \textbf{Lemma Statement}: \texttt{lemma m\_is\_odd (sol : ErdosMoserSolution) : Odd sol.m}
    \item \textbf{Proof Strategy}: \texttt{by\_contra} assuming \texttt{Even sol.m}, case split on parity of \texttt{q = sol.m / 2}, and analysis of 2-adic valuation using \texttt{padic\_valNat}.
\end{itemize}

\end{document}
